\documentclass[12pt]{article}
\usepackage{amsmath, amssymb}
\usepackage{geometry}
\usepackage{enumitem}
\usepackage{fancyhdr}
\usepackage{mdframed}
\usepackage[T1]{fontenc}
\usepackage{lmodern}

\geometry{top=2.5cm, bottom=2.5cm, left=2.5cm, right=2.5cm}
\setlength{\headheight}{15pt}
\pagestyle{fancy}
\fancyhf{}
\rhead{Infi 2 --- HIT (21152)}
\lhead{Changing the Order of Integration}
\cfoot{\thepage}
\renewcommand{\headrulewidth}{0.4pt}

\title{\textbf{Worksheet: Changing the Order of Integration}\\[4pt]
\large Double Integrals --- Practice Problems}
\author{Course 21152 --- Infinitesimal Calculus 2, HIT}
\date{}

\begin{document}
\maketitle
\thispagestyle{fancy}

\begin{mdframed}
\textbf{Procedure.}
\begin{enumerate}[nosep]
  \item Sketch the domain $D$ described by the given limits.
  \item Identify all boundary curves and their intersection points.
  \item Fix the \emph{new} outer variable and read its full range from the sketch.
  \item For each fixed value of the outer variable, find the range of the inner variable from the boundary curves (express as functions of the outer variable).
  \item Rewrite the integral.
\end{enumerate}
\end{mdframed}

\bigskip

%--------------------------------------------------------------
\section*{Part A \quad Change the Order of Integration}
\textit{Source: Reznik, \emph{Problems in Higher Mathematics}, pp.\ 95--96.}\\[4pt]
Rewrite each integral as an equivalent iterated integral with the \textbf{opposite order of integration}.
Note: the result may require splitting into two or three sub-integrals.

\begin{enumerate}[label=\textbf{A\arabic*.}, leftmargin=*, itemsep=18pt]

\item % Reznik 12
\[
  \int_0^{4} \int_{y}^{10-y} f(x,y)\,dx\,dy
\]
\textit{Domain bounded by $y = x$ and $y = 10-x$, with $y \ge 0$.}

\item % Reznik 13
\[
  \int_1^{3} \int_{x/3}^{2x} f(x,y)\,dy\,dx
\]
\textit{Domain bounded by $y = 2x$ and $y = x/3$, between $x=1$ and $x=3$.}

\item % Reznik 14
\[
  \int_0^{3} \int_0^{\sqrt{25-x^2}} f(x,y)\,dy\,dx
\]
\textit{Domain: region inside the circle $x^2+y^2=25$, first quadrant, with $x \le 3$.}

\item % Reznik 15
\[
  \int_{-1}^{2} \int_{x^2}^{x+2} f(x,y)\,dy\,dx
\]
\textit{Domain bounded by the parabola $y = x^2$ and the line $y = x+2$.}

\item % Reznik 16 (has parameter a)
($a > 0$)\;
\[
  \int_0^{2a} \int_{\sqrt{2ax-x^2}}^{\sqrt{4ax}} f(x,y)\,dy\,dx
\]
\textit{Hint: $y = \sqrt{2ax - x^2}$ rewrites as $(x-a)^2 + y^2 = a^2$ (upper semicircle of radius $a$ centred at $(a,0)$); $y = \sqrt{4ax}$ is a parabola with vertex at the origin.}

\item % Reznik 17
\[
  \int_0^{1} \int_{-\sqrt{1-y^2}}^{1-y} f(x,y)\,dx\,dy
\]
\textit{Domain: inside the unit circle $x^2+y^2=1$ to the left of the line $x = 1-y$, with $y \ge 0$.}

\item % Reznik 18
\[
  \int_0^{1} \int_{y^2/2}^{\sqrt{3-y^2}} f(x,y)\,dx\,dy
\]
\textit{Domain: to the right of $x = y^2/2$ (parabola) and inside the circle $x^2+y^2=3$, with $y \ge 0$.}

\end{enumerate}

\bigskip
%--------------------------------------------------------------
\section*{Part B \quad Reverse the Order to Evaluate}
\textit{Source: Hughes-Hallett et al., \emph{Calculus: Single and Multivariable}, 6th ed., \S16.2, problems 33--37.}\\[4pt]
Each integral \textbf{cannot} be evaluated as written (the inner antiderivative has no closed form).
Reverse the order of integration, then evaluate.

\begin{enumerate}[label=\textbf{B\arabic*.}, leftmargin=*, itemsep=18pt]

\item % HH 33
\[
  \int_0^1 \int_y^1 e^{x^2}\,dx\,dy
\]

\item % HH 34
\[
  \int_0^1 \int_y^1 \sin(x^2)\,dx\,dy
\]

\item % HH 35
\[
  \int_0^1 \int_{\sqrt{y}}^1 \sqrt{2+x^3}\,dx\,dy
\]

\item % HH 36
\[
  \int_0^3 \int_{y^2}^{9} y\sin(x^2)\,dx\,dy
\]

\item % HH 37
\[
  \int_0^1 \int_{e^y}^{e} \frac{x}{\ln x}\,dx\,dy
\]

\end{enumerate}

\bigskip
%--------------------------------------------------------------
\section*{Part C \quad Additional Problems}
\textit{Standard problems --- reverse the order and evaluate.}

\begin{enumerate}[label=\textbf{C\arabic*.}, leftmargin=*, itemsep=18pt]

\item
\[
  \int_0^1 \int_x^1 \sin(y^2)\,dy\,dx
\]
\textit{Hint: after reversing, the inner integral reduces to a single factor.}

\item
\[
  \int_0^1 \int_{\sqrt{x}}^1 e^{y^3}\,dy\,dx
\]

\item
\[
  \int_0^{\pi} \int_y^{\pi} \frac{\sin x}{x}\,dx\,dy
\]
\textit{Classic problem --- the reversed form evaluates to a familiar integral.}

\end{enumerate}

\end{document}
